% !TEX root = main.tex
% !TEX program = pdflatex
\documentclass[a4paper, 14pt]{extarticle} % 14pt вместо 12pt
\usepackage[T2A]{fontenc}
\usepackage[utf8]{inputenc}
\usepackage[russian]{babel}
\usepackage{mathptmx} % Times New Roman шрифт
\usepackage{geometry}
\usepackage{tabularx}
\usepackage{multirow}
\usepackage{booktabs}
\usepackage{amsmath}
\usepackage{graphicx}
\usepackage{float}
\usepackage{setspace} % Для межстрочного интервала
\usepackage{indentfirst} % Отступ первой строки абзаца
\usepackage{enumitem}
\usepackage{pdflscape}
\usepackage{rotating}
\usepackage{longtable}
\usepackage{booktabs}
\usepackage{array}

% Настройка страницы по ГОСТ
\geometry{
    left=30mm,    % 3 см
    right=15mm,   % 1,5 см
    top=20mm,     % 2 см
    bottom=20mm,  % 2 см
}

% Межстрочный интервал 1,5
\onehalfspacing

% Отступ первой строки абзаца 1,25 см
\setlength{\parindent}{1.25cm}

% Настройка заголовков
\usepackage{titlesec}

% Формат заголовков разделов (section)
\titleformat{\section}
    {\normalfont\Large\bfseries}  % размер: Large, шрифт: жирный
    {\thesection}                  % номер раздела
    {1em}                          % расстояние между номером и названием
    {}                             % код после заголовка

% Формат подразделов (subsection)
\titleformat{\subsection}
    {\normalfont\large\bfseries}   % размер: large, шрифт: жирный
    {\thesubsection}
    {1em}
    {}

% Формат под-подразделов (subsubsection)
\titleformat{\subsubsection}
    {\normalfont\normalsize\bfseries} % размер: normal, шрифт: жирный
    {\thesubsubsection}
    {1em}
    {}

% Установка отступов до и после заголовков
\titlespacing*{\section}{0pt}{12pt}{6pt}      % {влево}{до}{после}
\titlespacing*{\subsection}{0pt}{10pt}{4pt}
\titlespacing*{\subsubsection}{0pt}{8pt}{2pt}

% Убираем отступы в списках
\usepackage{enumitem}
\setlist{nosep, leftmargin=1.25cm}

% Для не нумерованных заголовков тоже
\titleformat{name=\section,numberless}
    {\normalfont\Large\bfseries}
    {}
    {0pt}
    {}

\titleformat{name=\subsection,numberless}
    {\normalfont\large\bfseries}
    {}
    {0pt}
    {}

% Настройка гиперссылок (рекомендуется подключать последним)
\usepackage{hyperref}
\hypersetup{
    colorlinks=true,
    linkcolor=blue, % Цвет ссылок
    urlcolor=blue
}


\usepackage[russian]{cleveref}
